\chapter{अभिक्रिया क्रियाविधियाँ: वक्र तीर}\label{ch:g12:curly-arrows}

\omterm{def:g8:balanced-equations:equation}{संतुलित समीकरण} किसी फ़िल्म की पहली और अंतिम तस्वीरों जैसा है: वह
पहले के \omterm{def:g7:chemical-reactions:reactant}{अभिकारक} और बाद के उत्पाद दिखाता है, और बीच का कुछ भी
नहीं। फिर भी \omterm{def:g7:atoms-and-molecules:molecule}{अणु} बस \omterm{def:g7:atoms-and-molecules:atom}{परमाणुओं} की अदला-बदली नहीं करते। आबंध एक-एक करके
टूटते और बनते हैं, जैसे-जैसे \omterm{def:g9:inside-the-atom:nucleus}{इलेक्ट्रॉन} युग्म एक जगह से दूसरी जगह जाते हैं,
ऐसी अल्पजीवी स्पीशीज़ से होकर जो समीकरण में कभी दिखाई नहीं देतीं। रसायनज्ञ
इन गतियों को मुड़े हुए तीरों से दर्शाते हैं। कुछ नियमों के साथ ये तीर
समझाते हैं कि कोई अभिक्रिया वहीं क्यों होती है जहाँ होती है, और उन अभिक्रियाओं का
पूर्वानुमान लगाते हैं जो पहले कभी देखी नहीं गईं।

\begin{recall}
\omterm{def:g11:polarity-and-cohesion:electronegativity}{विद्युत-ऋणात्मकता}, \omterm{def:g11:polarity-and-cohesion:polar-bond}{ध्रुवीय आबंध} और \omterm{def:g11:polarity-and-cohesion:polar-bond}{आंशिक आवेश}
(\cref{ch:g11:polarity-and-cohesion})। \omterm{def:g10:lewis-and-shape:lewis-structure}{लुइस संरचनाएँ}: \omterm{def:g10:lewis-and-shape:lone-pair}{आबंधी युग्म} और
\omterm{def:g10:lewis-and-shape:lone-pair}{एकाकी युग्म} (\cref{ch:g10:lewis-and-shape})। \omterm{def:g11:functional-groups:functional-group}{क्रियात्मक समूह} और
परिवार (\cref{ch:g11:functional-groups})।
\end{recall}

\section{इलेक्ट्रॉन-समृद्ध और इलेक्ट्रॉन-न्यून स्थल}

\begin{definition}[इलेक्ट्रॉन-दाता और इलेक्ट्रॉन-ग्राही स्थल]\label{def:g12:curly-arrows:site}
किसी \omterm{def:g7:atoms-and-molecules:molecule}{अणु} या \omterm{def:g9:ions:ion}{आयन} का \emph{इलेक्ट्रॉन-दाता स्थल}\index{इलेक्ट्रॉन-दाता स्थल}
\omterm{def:g9:inside-the-atom:nucleus}{इलेक्ट्रॉनों} से समृद्ध एक स्थान है, जो नया आबंध बनाने के लिए एक युग्म दे
सकता है: \omterm{def:g10:lewis-and-shape:lone-pair}{एकाकी युग्म}, \omterm{def:g10:lewis-and-shape:multiple-bond}{द्विआबंध}, ऋणात्मक आवेश या
\omterm{def:g11:polarity-and-cohesion:polar-bond}{आंशिक आवेश} $\delta^-$ वाला \omterm{def:g7:atoms-and-molecules:atom}{परमाणु}। \emph{इलेक्ट्रॉन-ग्राही
स्थल}\index{इलेक्ट्रॉन-ग्राही स्थल} \omterm{def:g9:inside-the-atom:nucleus}{इलेक्ट्रॉनों} में न्यून एक स्थान है, जो एक युग्म
ग्रहण कर सकता है: धनात्मक आवेश या \omterm{def:g11:polarity-and-cohesion:polar-bond}{आंशिक आवेश}
$\delta^+$ वाला \omterm{def:g7:atoms-and-molecules:atom}{परमाणु}।
\end{definition}

\begin{example}[दो अणुओं में स्थल]\label{ex:g12:curly-arrows:sites}
ब्रोमोएथेन, \ce{CH3-CH2-Br}, में ब्रोमीन कार्बन से अधिक विद्युत-ऋणात्मक है:
उससे आबंधित कार्बन पर $\delta^+$ है (ग्राही स्थल), और
ब्रोमीन पर $\delta^-$ और तीन \omterm{def:g10:lewis-and-shape:lone-pair}{एकाकी युग्म}। प्रोपेनोन में
\ce{C=O} समूह के ऑक्सीजन पर $\delta^-$ और दो \omterm{def:g10:lewis-and-shape:lone-pair}{एकाकी युग्म} हैं (दाता स्थल), उसके
कार्बन पर $\delta^+$ (ग्राही स्थल)। हाइड्रॉक्साइड \omterm{def:g9:ions:ion}{आयन} \ce{OH-}, अपने
ऋणात्मक आवेश और तीन \omterm{def:g10:lewis-and-shape:lone-pair}{एकाकी युग्मों} के साथ, एक प्रबल दाता है।
\end{example}

\begin{omfigure}
\setchemfig{atom sep=2.2em}
\begin{tabular}{cc}
\chemfig{H_3C-C(-[2]H)(-[6]H)-\charge{90:2pt=\:,0=\:,270:2pt=\:}{Br}} &
\chemfig{H_3C-C(-[6]CH_3)=[2]\charge{0=\:,180=\:}{O}} \\[2pt]
\small \ce{Br} से आबंधित कार्बन पर $\delta^+$, \ce{Br} पर $\delta^-$ &
\small \ce{C=O} के कार्बन पर $\delta^+$, \ce{O} पर $\delta^-$ \\
\small ब्रोमोएथेन & \small प्रोपेनोन
\end{tabular}

{\small इलेक्ट्रॉन-न्यून कार्बन \omterm{def:g7:atoms-and-molecules:atom}{परमाणु} ($\delta^+$, ग्राही स्थल), इलेक्ट्रॉन-समृद्ध,
विद्युत-ऋणात्मक \omterm{def:g7:atoms-and-molecules:atom}{परमाणुओं} ($\delta^-$, \omterm{def:g10:lewis-and-shape:lone-pair}{एकाकी युग्मों} के साथ: दाता स्थल) के बगल में।}
\end{omfigure}

\section{वक्र तीर}

\begin{definition}[वक्र तीर]\label{def:g12:curly-arrows:curly-arrow}
\emph{वक्र तीर}\index{वक्र तीर} किसी अभिक्रिया के एक चरण के दौरान
\omterm{def:g9:inside-the-atom:nucleus}{इलेक्ट्रॉनों} के एक युग्म की गति दिखाता है। वह गति करने वाले युग्म से,
किसी \omterm{def:g10:lewis-and-shape:lone-pair}{एकाकी युग्म} या आबंध से, शुरू होता है और वहाँ समाप्त होता है जहाँ युग्म जाता है: किसी \omterm{def:g7:atoms-and-molecules:atom}{परमाणु} पर,
जहाँ वह नया आबंध बनाता है, या उस \omterm{def:g7:atoms-and-molecules:atom}{परमाणु} पर जो टूटने वाले आबंध
का युग्म ले लेता है।
\end{definition}

\begin{method}[वक्र तीर बनाना]\label{met:g12:curly-arrows:drawing}
\begin{enumerate}
  \item दाता स्थल और ग्राही स्थल खोजिए।
  \item दाता युग्म (\omterm{def:g10:lewis-and-shape:lone-pair}{एकाकी युग्म} या आबंध) से ग्राही \omterm{def:g7:atoms-and-molecules:atom}{परमाणु} तक एक तीर
        बनाइए: वहाँ एक नया आबंध बनता है।
  \item यदि तब ग्राही \omterm{def:g7:atoms-and-molecules:atom}{परमाणु} पर बहुत अधिक \omterm{def:g9:inside-the-atom:nucleus}{इलेक्ट्रॉन} हो जाएँ (अष्टक से
        अधिक, या हाइड्रोजन के लिए दो से अधिक), तो उसके किसी आबंध को तोड़ने वाला
        दूसरा तीर बनाइए, जिसका युग्म अधिक
        विद्युत-ऋणात्मक \omterm{def:g7:atoms-and-molecules:atom}{परमाणु} पर जाता है।
  \item चरण के बाद आवेश जाँचिए: \omterm{def:g10:lewis-and-shape:lone-pair}{एकाकी युग्म} देने वाला \omterm{def:g7:atoms-and-molecules:atom}{परमाणु}
        एक अधिक धनात्मक हो जाता है, \omterm{def:g10:lewis-and-shape:lone-pair}{एकाकी युग्म} के रूप में युग्म ग्रहण करने वाला \omterm{def:g7:atoms-and-molecules:atom}{परमाणु}
        एक अधिक ऋणात्मक; कुल आवेश नहीं बदलता।
\end{enumerate}
\end{method}

\begin{proposition}[हर चरण में आवेश संरक्षित रहता है]\label{prop:g12:curly-arrows:charge}
किसी क्रियाविधि के हर चरण से पहले और बाद में स्पीशीज़ का कुल विद्युत आवेश
एक ही रहता है।
\end{proposition}

\begin{proof}
\omterm{def:g12:curly-arrows:curly-arrow}{वक्र तीर} केवल पहले से मौजूद \omterm{def:g9:inside-the-atom:nucleus}{इलेक्ट्रॉनों} को खिसकाता है; कोई \omterm{def:g9:inside-the-atom:nucleus}{इलेक्ट्रॉन}
न बनता है न नष्ट होता है, और \omterm{def:g9:inside-the-atom:nucleus}{नाभिक} नहीं बदलते।
\end{proof}

\begin{example}[पहला तीर]\label{ex:g12:curly-arrows:first}
एक हाइड्रोजन \omterm{def:g9:ions:ion}{आयन} \ce{H+} एक हाइड्रॉक्साइड \omterm{def:g9:ions:ion}{आयन} \ce{OH-} से मिलता है। ऑक्सीजन का
एक \omterm{def:g10:lewis-and-shape:lone-pair}{एकाकी युग्म} उस हाइड्रोजन \omterm{def:g9:ions:ion}{आयन} के साथ आबंध बनाता है, जिसके पास कोई \omterm{def:g9:inside-the-atom:nucleus}{इलेक्ट्रॉन} नहीं:
\ce{H+ + OH- -> H2O}। एक तीर, ऑक्सीजन के \omterm{def:g10:lewis-and-shape:lone-pair}{एकाकी युग्म} से
\ce{H+} तक; आवेश $+1$ और $-1$ पानी के \omterm{def:g7:atoms-and-molecules:molecule}{अणु} में 0 हो जाते हैं।
\end{example}

\begin{omfigure}
\setchemfig{atom sep=2.1em}
\chemfig{@{hp}H^{+}}\qquad\raisebox{0.2ex}{$+$}\qquad
\chemfig{H-@{ox}\charge{90=\:,0=\:,270=\:}{O}^{-}}
\qquad\raisebox{0.2ex}{$\longrightarrow$}\qquad
\chemfig{H-\charge{90=\:,270=\:}{O}-H}
\chemmove{
  \draw[->, thick, omProp, shorten <=4pt, shorten >=3pt]
    ($(ox)+(0,0.25)$) .. controls +(110:9mm) and +(60:9mm) .. (hp);}

{\small सबसे सरल चरण: हाइड्रॉक्साइड \omterm{def:g9:ions:ion}{आयन} का एक \omterm{def:g10:lewis-and-shape:lone-pair}{एकाकी युग्म} खिसककर
हाइड्रोजन \omterm{def:g9:ions:ion}{आयन} के साथ नया \ce{O-H} आबंध बनाता है।}
\end{omfigure}

\section{मूल चरण और मध्यवर्ती}

\begin{definition}[क्रियाविधि, मूल चरण, मध्यवर्ती]\label{def:g12:curly-arrows:mechanism}
किसी अभिक्रिया की \emph{अभिक्रिया क्रियाविधि}\index{अभिक्रिया क्रियाविधि}
चरणों की वह श्रृंखला है जिससे \omterm{def:g7:chemical-reactions:reactant}{अभिकारक} उत्पाद बनते हैं। हर
\emph{मूल चरण}\index{मूल चरण} एक अकेली घटना है जिसमें
\omterm{def:g9:inside-the-atom:nucleus}{इलेक्ट्रॉनों} के कुछ युग्म एक साथ गति करते हैं। एक चरण में बनने और बाद के किसी चरण में
खपने वाली स्पीशीज़ एक \emph{अभिक्रिया मध्यवर्ती}\index{अभिक्रिया मध्यवर्ती} है:
वह समग्र समीकरण में दिखाई नहीं देती।
\end{definition}


\begin{example}[एक कार्बधनायन मध्यवर्ती]\label{ex:g12:curly-arrows:carbocation}
जब हाइड्रोजन ब्रोमाइड एथीन से जुड़ता है, तो अभिक्रिया दो चरणों में होती है।
पहले चरण में एथीन का \omterm{def:g10:lewis-and-shape:multiple-bond}{द्विआबंध}
\ce{HBr} के हाइड्रोजन को एक युग्म देता है, और \ce{H-Br} युग्म ब्रोमीन पर जाता है, जो
\ce{Br-} के रूप में निकल जाता है; जिस कार्बन ने \omterm{def:g10:lewis-and-shape:multiple-bond}{द्विआबंध} में अपना हिस्सा खोया, उसके पास
केवल छह \omterm{def:g9:inside-the-atom:nucleus}{इलेक्ट्रॉन} और एक धनात्मक आवेश रह जाता है: एक कार्बधनायन,
\ce{CH3-CH2+}। दूसरे चरण में \ce{Br-} का एक \omterm{def:g10:lewis-and-shape:lone-pair}{एकाकी युग्म} उस कार्बन के साथ
आबंध बनाता है। कार्बधनायन एक मध्यवर्ती है।
\end{example}

\section{अभिक्रियाओं की तीन श्रेणियाँ}

\begin{definition}[प्रतिस्थापन, योगज, विलोपन]\label{def:g12:curly-arrows:categories}
\emph{प्रतिस्थापन}\index{प्रतिस्थापन} में किसी \omterm{def:g7:atoms-and-molecules:molecule}{अणु} का कोई \omterm{def:g7:atoms-and-molecules:atom}{परमाणु} या समूह
दूसरे से बदल जाता है। \emph{योगज अभिक्रिया}\index{योगज अभिक्रिया} में
दो \omterm{def:g7:atoms-and-molecules:molecule}{अणु} मिलकर एक हो जाते हैं, और एक बहु-आबंध एकल आबंध बन जाता है।
\emph{विलोपन}\index{विलोपन} में किसी \omterm{def:g7:atoms-and-molecules:molecule}{अणु} से एक छोटा \omterm{def:g7:atoms-and-molecules:molecule}{अणु} (अक्सर पानी)
हटा दिया जाता है, और एक एकल आबंध बहु-आबंध बन जाता है।
\end{definition}

\begin{omfigure}
\setchemfig{atom sep=2em}
\small
\textbf{\omterm{def:g12:curly-arrows:categories}{प्रतिस्थापन}} (एक चरण): \ce{CH3-CH2-Br + OH- -> CH3-CH2-OH + Br-}

\medskip
\chemfig{H-@{o1}\charge{90=\:,0=\:,270=\:}{O}^{-}}\qquad
\chemfig{H_3C-@{c1}C(-[2]H)(-[6]H)-[@{b1}]@{br1}\charge{90=\:,0=\:,270=\:}{Br}}
\qquad$\longrightarrow$\qquad
\chemfig{H_3C-C(-[2]H)(-[6]H)-O-H}\quad$+$\quad\ce{Br-}
\chemmove{
  \draw[->, thick, omProp, shorten <=4pt, shorten >=3pt]
    ($(o1)+(0.15,0.2)$) .. controls +(60:9mm) and +(120:9mm) .. (c1);
  \draw[->, thick, omDef, shorten <=2pt, shorten >=4pt]
    (b1) .. controls +(-90:7mm) and +(-90:7mm) .. (br1);}

\vspace{6mm}
\textbf{योगज} (दो चरण): \ce{CH2=CH2 + HBr -> CH3-CH2-Br}

\vspace{9mm}
\chemfig{H_2C=[@{pi}]CH_2}\quad$+$\quad
\chemfig{@{h2}H-[@{b2}]@{br2}Br}
\qquad$\longrightarrow$\qquad
\chemfig{H_3C-@{cp}\charge{90:2pt=$\scriptstyle +$}{C}H_2}\quad$+$\quad
\chemfig{@{brm}\charge{90=\:,180=\:,270=\:,0=\:}{Br}^{-}}
\qquad$\longrightarrow$\qquad \ce{CH3-CH2-Br}
\chemmove{
  \draw[->, thick, omProp, shorten <=2pt, shorten >=3pt]
    (pi) .. controls +(90:10mm) and +(110:9mm) .. (h2);
  \draw[->, thick, omDef, shorten <=2pt, shorten >=4pt]
    (b2) .. controls +(-90:7mm) and +(-90:7mm) .. (br2);
  \draw[->, thick, omProp, shorten <=4pt, shorten >=3pt]
    ($(brm)+(0,-0.25)$) .. controls +(-120:8mm) and +(-60:8mm) .. (cp);}

\vspace{6mm}
\textbf{\omterm{def:g12:curly-arrows:categories}{विलोपन}} (एथेनॉल का अम्लीय निर्जलीकरण, तीन चरण):
\ce{CH3-CH2-OH -> CH2=CH2 + H2O}

\medskip
\begin{tabular}{@{}l@{}}
1.\ ऑक्सीजन एक हाइड्रोजन \omterm{def:g9:ions:ion}{आयन} लेता है: \ce{CH3-CH2-OH + H+ -> CH3-CH2-OH2+};\\
2.\ \ce{C-O} युग्म पानी के साथ निकल जाता है: \ce{CH3-CH2-OH2+ -> CH3-CH2+ + H2O};\\
3.\ पड़ोसी \ce{C-H} युग्म \omterm{def:g10:lewis-and-shape:multiple-bond}{द्विआबंध} बन जाता है, और
    \ce{H+} निकल जाता है:\\
\phantom{3.\ }\ce{CH3-CH2+ -> CH2=CH2 + H+}।
\end{tabular}

\medskip
{\small तीन क्रियाविधियाँ। नारंगी तीर: दाता युग्म एक नया आबंध बनाता है।
नीले तीर: एक आबंध टूटता है, उसका युग्म अधिक विद्युत-ऋणात्मक
\omterm{def:g7:atoms-and-molecules:atom}{परमाणु} पर जाता है। \omterm{def:g12:curly-arrows:categories}{विलोपन} में हाइड्रोजन \omterm{def:g9:ions:ion}{आयन} अंत में वापस मिल जाता है।}
\end{omfigure}


\begin{remark}[कंकाल या समूह बदलना]\label{rem:g12:curly-arrows:skeleton}
कुछ अभिक्रियाएँ केवल \omterm{def:g11:functional-groups:functional-group}{क्रियात्मक समूह} बदलती हैं और कार्बन
कंकाल रखती हैं: ब्रोमोएथेन का \omterm{def:g12:curly-arrows:categories}{प्रतिस्थापन} एक \omterm{def:g11:functional-groups:halogenoalkane}{हैलोऐल्केन} को
\omterm{def:g11:functional-groups:alcohol}{ऐल्कोहॉल} में बदलता है, पहले दो कार्बन, बाद में भी दो। दूसरी अभिक्रियाएँ
कंकाल बनाती या काटती हैं: नया कार्बन--कार्बन आबंध बनाना श्रृंखला को लंबा करता है, जो
छोटे \omterm{def:g7:atoms-and-molecules:molecule}{अणुओं} से बड़ा \omterm{def:g7:atoms-and-molecules:molecule}{अणु} बनाने का मुख्य चरण है। \omterm{def:g10:synthesis-yield:synthesis}{संश्लेषण} की योजना बनाने का
अर्थ है दोनों प्रकार की अभिक्रियाओं को सही क्रम में चुनना।
\end{remark}

\section{अभ्यास}

\begin{exercise}[$\star$]\label{exo:g12:curly-arrows:1}
हर स्पीशीज़ में एक दाता स्थल और, यदि हो, तो एक ग्राही
स्थल बताइए: \ce{OH-}; \ce{H+}; \ce{CH2=CH2}; \ce{CH3-Cl}; \ce{H2O}।
\end{exercise}

\begin{exercise}[$\star$]\label{exo:g12:curly-arrows:2}
\omterm{def:g12:curly-arrows:categories}{प्रतिस्थापन}, योगज या \omterm{def:g12:curly-arrows:categories}{विलोपन}?
(a) \ce{CH3-CH2-Cl + OH- -> CH3-CH2-OH + Cl-};
(b) \ce{CH2=CH2 + H2O -> CH3-CH2-OH};
(c) \ce{CH3-CH2-OH -> CH2=CH2 + H2O};
(d) \ce{CH2=CH2 + Br2 -> CH2Br-CH2Br}।
\end{exercise}

\begin{exercise}[$\star$]\label{exo:g12:curly-arrows:3}
\omterm{def:g12:curly-arrows:curly-arrow}{वक्र तीर} कहाँ से शुरू होता है, और कहाँ समाप्त होता है? किसी ऑक्सीजन \omterm{def:g7:atoms-and-molecules:atom}{परमाणु} के \omterm{def:g10:lewis-and-shape:lone-pair}{एकाकी युग्म}
से किसी कार्बन \omterm{def:g7:atoms-and-molecules:atom}{परमाणु} तक के \omterm{def:g12:curly-arrows:curly-arrow}{वक्र तीर} का क्या अर्थ है?
\end{exercise}

\begin{exercise}[$\star$]\label{exo:g12:curly-arrows:4}
अभिक्रिया \ce{H+ + NH3 -> NH4+} में कौन-सी स्पीशीज़ युग्म देती है?
तीर बनाइए। उत्पाद का आवेश क्या है?
\end{exercise}

\begin{exercise}[$\star$]\label{exo:g12:curly-arrows:5}
\omterm{def:g12:curly-arrows:mechanism}{अभिक्रिया मध्यवर्ती} क्या है? एथीन से
\ce{HBr} के योग के मध्यवर्ती का नाम बताइए।
\end{exercise}

\begin{exercise}[$\star\star$]\label{exo:g12:curly-arrows:6}
\omterm{def:g12:curly-arrows:categories}{प्रतिस्थापन} \ce{CH3-Br + OH- -> CH3-OH + Br-} (एक चरण) के \omterm{def:g12:curly-arrows:curly-arrow}{वक्र तीर}
बनाइए, और आवेश जाँचिए।
\end{exercise}

\begin{exercise}[$\star\star$]\label{exo:g12:curly-arrows:7}
एथीन से \ce{HBr} के योग के पहले चरण में
\ce{H-Br} युग्म हाइड्रोजन के बजाय ब्रोमीन पर क्यों जाता है? चरण के हर
उत्पाद पर कौन-सा आवेश होता है?
\end{exercise}

\begin{exercise}[$\star\star$]\label{exo:g12:curly-arrows:8}
एथेनॉल के निर्जलीकरण में तीनों चरण लिखिए और हर एक के लिए बताइए
कि कौन-सी स्पीशीज़ दाता है और कौन-सी ग्राही। हाइड्रोजन
\omterm{def:g9:ions:ion}{आयन} समग्र समीकरण में क्यों नहीं लिखा जाता?
\end{exercise}

\begin{exercise}[$\star\star$]\label{exo:g12:curly-arrows:9}
चरण \ce{CH3-CH2+ + H2O -> CH3-CH2-OH2+} पानी में एथीन से हाइड्रोजन \omterm{def:g9:ions:ion}{आयन}
के योग के बाद आता है। तीर बनाइए। \ce{H+} के एक और निष्कासन के बाद क्या
बनता है? समग्र समीकरण लिखिए और
उसकी श्रेणी बताइए।
\end{exercise}

\begin{exercise}[$\star\star$]\label{exo:g12:curly-arrows:10}
ब्रोमोएथेन के \omterm{def:g12:curly-arrows:categories}{प्रतिस्थापन} की क्रियाविधि पढ़िए: कितने युग्म
गति करते हैं? कौन-सा आबंध बनता है, कौन-सा टूटता है? क्या कोई मध्यवर्ती है?
\end{exercise}

\begin{exercise}[$\star\star$]\label{exo:g12:curly-arrows:11}
% ledger: en:C, en:Cl, en:Br
क्लोरोमेथेन हाइड्रॉक्साइड \omterm{def:g9:ions:ion}{आयन} से ब्रोमोमेथेन की तुलना में अधिक धीरे
अभिक्रिया करता है। \omterm{def:g11:polarity-and-cohesion:electronegativity}{विद्युत-ऋणात्मकताओं} (C 2.55, Cl 3.16, Br 2.96) की मदद से
समझाइए कि कौन-सा कार्बन अधिक $\delta^+$ है। क्या इससे
चालें समझ में आती हैं? (यह भी सोचिए कि निकलने वाला \omterm{def:g9:ions:ion}{आयन} युग्म को कितनी आसानी से लेता है।)
\end{exercise}

\begin{exercise}[$\star\star\star$]\label{exo:g12:curly-arrows:12}
हाइड्रोजन क्लोराइड प्रोपीन, \ce{CH2=CH-CH3}, से जुड़ता है। पहले चरण में
हाइड्रोजन \omterm{def:g9:ions:ion}{आयन} \omterm{def:g10:lewis-and-shape:multiple-bond}{द्विआबंध} के किसी भी कार्बन से आबंधित हो सकता है, जिससे
दो संभावित कार्बधनायन बनते हैं। दोनों बनाइए, दोनों संभावित उत्पाद भी, और उनके
नाम बताइए। (कौन-सा अधिक बनता है, इसका अध्ययन विश्वविद्यालय के प्रथम वर्ष के खंड में है।)
\end{exercise}

\begin{exercise}[$\star\star\star$]\label{exo:g12:curly-arrows:13}
एक छात्र ब्रोमोएथेन के \omterm{def:g12:curly-arrows:categories}{प्रतिस्थापन} का पहला चरण ब्रोमीन \omterm{def:g7:atoms-and-molecules:atom}{परमाणु} से
\ce{OH-} के ऑक्सीजन तक एक तीर के साथ बनाता है। दोनों
ग़लतियाँ समझाइए।
\end{exercise}

\begin{exercise}[$\star\star\star$]\label{exo:g12:curly-arrows:14}
\omterm{def:g11:functional-groups:carboxylic-acid}{कार्बोक्सिलिक अम्ल} और \omterm{def:g11:functional-groups:alcohol}{ऐल्कोहॉल} से \omterm{def:g11:functional-groups:carboxylic-acid}{एस्टर} बनने में
\omterm{def:g11:functional-groups:alcohol}{ऐल्कोहॉल} का ऑक्सीजन \ce{C=O} समूह के कार्बन पर आक्रमण करता है। \omterm{def:g11:polarity-and-cohesion:polar-bond}{आंशिक आवेशों}
की मदद से समझाइए कि वह कार्बन ग्राही स्थल और वह
ऑक्सीजन दाता स्थल क्यों है। समग्र रूप से कौन-सा छोटा \omterm{def:g7:atoms-and-molecules:molecule}{अणु} निकलता है, और पूरी
अभिक्रिया किस श्रेणी की है?
\end{exercise}

\begin{exercise}[$\star\star\star$]\label{exo:g12:curly-arrows:15}
एथीन ब्रोमीन \ce{Br2}, एक \omterm{def:g11:polarity-and-cohesion:polar-molecule}{अध्रुवीय अणु}, से अभिक्रिया करता है। जैसे ही \ce{Br2}
\omterm{def:g7:atoms-and-molecules:molecule}{अणु} \omterm{def:g10:lewis-and-shape:multiple-bond}{द्विआबंध} के पास आता है, \ce{Br-Br} के \omterm{def:g9:inside-the-atom:nucleus}{इलेक्ट्रॉन} निकट वाले ब्रोमीन \omterm{def:g7:atoms-and-molecules:atom}{परमाणु} से
दूर धकेले जाते हैं। समझाइए कि इससे ग्राही स्थल क्यों बनता है, और
क्रियाविधि का पहला चरण प्रस्तावित कीजिए।
\end{exercise}

\section{समस्या: ब्रोमोएथेन बनाना}

\begin{problem}[{सप्ताहांत समस्या --- दस ग्राम एथीन से: कितना
ब्रोमोएथेन?}]
\label{pb:g12:curly-arrows:1}
% ledger: aw:C, aw:H, aw:Br, en:H, en:Br, en:C
ब्रोमोएथेन, जो अनेक \omterm{def:g10:synthesis-yield:synthesis}{संश्लेषणों} की प्रारंभिक सामग्री है, एथीन में
हाइड्रोजन ब्रोमाइड जोड़कर बनाया जाता है: \ce{CH2=CH2 + HBr -> CH3-CH2-Br}। एक रसायनज्ञ
\qty{10.0}{g} एथीन और हाइड्रोजन ब्रोमाइड की अधिकता इस्तेमाल करता है;
\omterm{def:g10:synthesis-yield:yield}{लब्धि} \qty{75}{\percent} है।

\medskip
\noindent\textbf{भाग I --- स्थल।}
\begin{enumerate}
  \item एथीन और हाइड्रोजन ब्रोमाइड की \omterm{def:g10:lewis-and-shape:lewis-structure}{लुइस संरचनाएँ} बनाइए।
  \item एथीन का दाता स्थल कहाँ है? क्यों?
  \item \ce{H-Br} आबंध का \omterm{def:g11:polarity-and-cohesion:electronegativity}{विद्युत-ऋणात्मकता} अंतर निकालिए
        (H 2.20, Br 2.96)। किस \omterm{def:g7:atoms-and-molecules:atom}{परमाणु} पर $\delta^+$ है?
  \item \ce{HBr} का कौन-सा \omterm{def:g7:atoms-and-molecules:atom}{परमाणु} ग्राही स्थल है?
\end{enumerate}

\medskip
\noindent\textbf{भाग II --- क्रियाविधि।}
\begin{enumerate}[resume]
  \item पहला चरण उसके दोनों \omterm{def:g12:curly-arrows:curly-arrow}{वक्र तीरों} के साथ बनाइए।
  \item बनने वाली दोनों स्पीशीज़ और उनके आवेश बताइए।
  \item पहले चरण में आवेश संरक्षण जाँचिए।
  \item दूसरा चरण उसके \omterm{def:g12:curly-arrows:curly-arrow}{वक्र तीर} के साथ बनाइए।
  \item मध्यवर्ती के धनात्मक कार्बन के चारों ओर कितने \omterm{def:g9:inside-the-atom:nucleus}{इलेक्ट्रॉन} हैं?
        वह इतना क्रियाशील क्यों है?
\end{enumerate}

\medskip
\noindent\textbf{भाग III --- किस प्रकार की अभिक्रिया?}
\begin{enumerate}[resume]
  \item क्रियाविधि का मध्यवर्ती क्या है?
  \item वह समग्र समीकरण में क्यों दिखाई नहीं देता?
  \item अभिक्रिया किस श्रेणी की है?
  \item क्या अभिक्रिया कार्बन कंकाल बदलती है, या केवल
        \omterm{def:g11:functional-groups:functional-group}{क्रियात्मक समूह}?
\end{enumerate}

\medskip
\noindent\textbf{भाग IV --- कितना उत्पाद?}
\begin{enumerate}[resume]
  \item एथीन और ब्रोमोएथेन के \omterm{def:g10:the-mole:molar-mass}{मोलर द्रव्यमान} निकालिए।
  \item एथीन की मात्रा निकालिए।
  \item \omterm{def:g10:reaction-progress-table:extent}{प्रगति सारणी} भरिए। कौन-सा \omterm{def:g7:chemical-reactions:reactant}{अभिकारक} सीमांत है?
  \item ब्रोमोएथेन का अधिकतम द्रव्यमान निकालिए।
  \item \qty{75}{\percent} की \omterm{def:g10:synthesis-yield:yield}{लब्धि} के साथ प्राप्त द्रव्यमान निकालिए।
  \item अंतिम उत्तर बताइए: रसायनज्ञ को ब्रोमोएथेन का कितना द्रव्यमान
        मिलता है?
\end{enumerate}
\end{problem}
